
% nonallelic genes
% localization signal / zipcode


\newglossaryentry{biovorl}{
    name={test},
    description={todo},
    type=defi
}

%% Noch wo rein

\newglossaryentry{crossingover}{
    name={Crossing-over / Rekombination},
    description={Austausch von Chromosomenstuecken zwischen homologen Chromosomen in der Meiose. Je weiter zwei Loci auseinanderliegen, desto haeufiger findet dazwischen ein Crossing-over statt -> hoehere Rekombinationsfrequenz.},
    type=dinge
}


%% Schon drin
\newglossaryentry{biparental}{
    name={Biparentale (extranukleaere) Vererbung},
    description={Organellengenome werden von beiden Elternteilen weitergegeben. Zwischen rein maternaler und rein paternaler Vererbung existieren viele Uebergangsformen.},
    type=defi
}

\newglossaryentry{extranukleaer}{
    name={Extranukläere Gene},
    description={Gene ausserhalb des Zellkerns (in Mitochondrium oder Chloroplast). Sie werden in demselben Organellenkompartiment transkribiert und translatiert, in dem sie liegen.},
    type=defi
}

\newglossaryentry{annotation}{
    name={Annotation},
    description={Das Zuweisen von Struktur und Funktion zu einer DNA-Sequenz: Auffinden von Genen, Exon-Intron-Struktur, kodierenden Bereichen und funktionelle Charakterisierung.},
    type=defi
}

\newglossaryentry{somatsegregation}{
    name={Somatische Segregation},
    description={Zufaellige Aufteilung unterschiedlicher Organellengenome auf Tochterzellen bei der Zellteilung (relevant bei Heteroplasmie, v.a. in Pflanzen), sodass verschiedene Gewebe unterschiedliche Organellengenotypen tragen koennen.},
    type=defi
}


\newglossaryentry{kopplungskarte}{
    name={Kopplungskarte (genetische Karte)},
    description={Zeigt Abstaende zwischen Loci in Einheiten, die auf Rekombinationshaeufigkeiten beruhen. Begrenzt durch die Notwendigkeit, Rekombinationen zwischen variablen Markern zu beobachten. Die Rekombinationshaeufigkeit ist proportional zum physikalischen Abstand der Loci.},
    type=defi
}

\newglossaryentry{kopplung}{
    name={Genetische Kopplung},
    description={Tendenz von Genen auf demselben Chromosom, gemeinsam an die Nachkommen vererbt zu werden, statt sich unabhaengig zu verteilen (Abweichung von Mendels Unabhaengigkeitsregel). Grundlage fuer genetische Karten.},
    type=defi
}

\newglossaryentry{contig}{
    name={Contig},
    description={Aus ueberlappenden Klonen bzw. Sequenz-Reads zusammengesetzter, zusammenhaengender ("contiguous") DNA-Bereich.},
    type=defi
}


\newglossaryentry{bac}{
    name={BAC (Bacterial Artificial Chromosome)},
    description={Auf dem F-Plasmid basierender Klonierungsvektor, der sehr grosse DNA-Fragmente (ca. 100--300 kb) stabil in E. coli aufnehmen kann. Grundlage fuer physikalische Karten.},
    type=defi
}

\newglossaryentry{altspleissen}{
    name={Alternatives Spleissen},
    description={Exons koennen unterschiedlich kombiniert werden, sodass aus einem Gen mehrere verschiedene mRNAs/Proteine entstehen. Ein Grund, warum "ein Gen = ein Protein" fuer Eukaryoten nicht zutrifft.},
    type=defi
}

\newglossaryentry{cwertparadox}{
    name={C-Wert-Paradoxon},
    description={Der C-Wert ist die haploide Genomgroesse (DNA-Menge) einer Art. Paradox: die Genomgroesse korreliert bei Eukaryoten nicht mit der Komplexitaet oder Genzahl des Organismus (grosse Genome enthalten viel nicht-kodierende/repetitive DNA).},
    type=defi
}

\newglossaryentry{ras}{
    name={RAS},
    description={todo},
    type=defi
}

\newglossaryentry{exchange_factor}{
    name={exchange factor},
    description={todo},
    type=defi
}

\newglossaryentry{wachstumsfaktor}{
    name={Wachstumsfaktor},
    description={todo},
    type=defi
}

\newglossaryentry{gi}{
    name={G1-Phase},
    description={todo},
    type=defi
}

\newglossaryentry{sphase}{
    name={S-Phase},
    description={todo},
    type=defi
}

\newglossaryentry{gii}{
    name={G2-Phase},
    description={todo},
    type=defi
}

\newglossaryentry{interphase}{
    name={Interphase},
    description={todo},
    type=defi
}

\newglossaryentry{telophase}{
    name={Telophase},
    description={die Chromatiden erreichen die entgegen-
    gesetzten Pole, Rückbildung der Kernhülle,
    Dekondensierung der Chromosomen,
    Nucleoli wieder sichtbar},
    type=defi
}


\newglossaryentry{prophase}{
    name={Prophase},
    description={Kondensierung des Chromatins,
    Chromosomen werden sichtbar,
    Spindelapparat bildet sich},
    type=defi
}

\newglossaryentry{prometaphase}{
    name={Prometaphase},
    description={Kernhülle zerfällt, Kinetochore,
    Spindelapparat vollständig},
    type=defi
}


\newglossaryentry{metaphase}{
    name={Metaphase},
    description={Bildung der Metaphasenplatte:
    Chromosomen ordnen sich in einer Ebene
    zwischen den Spindelpolen},
    type=defi
}


\newglossaryentry{anaphase}{
    name={Anaphase},
    description={gepaarte Kinetochore trennen sich,
    Chromatiden wandern zu den entgegengesetzten Polen},
    type=defi
}

\newglossaryentry{cytokinese}{
    name={Cytokinese},
    description={Trennung der Zelle in zwei Tochterzellen nach der Mitose},
    type=defi
}
